% © 2026 Ing. David Jaroš — CC BY-NC-ND 4.0
%
% This work is licensed under a Creative Commons Attribution-NonCommercial-NoDerivatives
% 4.0 International License (CC BY-NC-ND 4.0).
%
% License History: Earlier drafts (up to v0.3) were released under CC BY 4.0.
% From v0.4 onward, all material is released under CC BY-NC-ND 4.0 to protect
% the integrity of the theoretical work during ongoing academic development.

% UBT-AI-PROVENANCE-BEGIN
% schema: ubt-ai-provenance/v1
% tier: C_working
% ai_assistance: disclosed
% human_review: risk-based
% editorial_responsibility: Ing. David Jaroš
% policy: ../../AI_PROVENANCE.md
% notice: Working material; exhaustive human review is not claimed.
% UBT-AI-PROVENANCE-END

\documentclass[11pt,a4paper]{article}
\usepackage{amsmath,amssymb,amsthm}
\usepackage{booktabs}
\usepackage{geometry}
\geometry{margin=2.5cm}

\title{Gauge Kinetic Normalization in the UBT-to-QED Reduction}
\author{Ing.\ David Jaroš}
\date{2026-05-10}

\begin{document}
\maketitle

\section{Objective}
We track normalization factors entering
\begin{equation}
S_{\mathrm{eff}}\supset -\frac{1}{4e^2}\int d^4x\,\mathcal{F}_{\mu\nu}\mathcal{F}^{\mu\nu},
\end{equation}
starting from the canonical UBT gauge terms.

\section{Canonical normalization basis}
The electroweak gauge terms in canonical form are
\begin{equation}
\mathcal{L}_{\mathrm{gauge}}=-\frac14 W^i_{\mu\nu}W^{i\mu\nu}-\frac14 B_{\mu\nu}B^{\mu\nu}.
\end{equation}
For non-abelian sectors, generator traces are normalized as
\begin{equation}
\mathrm{Tr}(T^aT^b)=\frac12\delta^{ab},\qquad
\mathrm{Tr}(\tau^i\tau^j)=\frac12\delta^{ij},
\end{equation}
which fixes relative normalization conventions but does not by itself fix the absolute
abelian coupling normalization.

\section{Factor bookkeeping}
\begin{center}
\begin{tabular}{@{}lll@{}}
\toprule
Source & Factor & Role \\
\midrule
Gauge kinetic definition & $1/4$ & Antisymmetric tensor normalization \\
Electromagnetic conversion & $\mathcal{A}_\mu=eA_\mu$ & Moves $e$ from matter vertex to kinetic prefactor \\
Fine-structure definition & $4\pi$ & $\alpha=e^2/(4\pi)$ in Heaviside-Lorentz units \\
Compact phase volume & $2\pi R_\psi$ & Appears after $\psi$ integration of higher-dimensional action \\
Zero-mode normalization & $(2\pi R_\psi)^{-1/2}$ & Cancels volume in canonically normalized 4D fields \\
EW mixing & $(\sin\theta_W,\cos\theta_W)$ & Orthogonal rotation, no extra Jacobian \\
\bottomrule
\end{tabular}
\end{center}

\section{Compactification and coupling normalization}
If one starts from a higher-dimensional kinetic term on
$\mathcal{M}_4\times S^1_\psi$ with circumference $2\pi R_\psi$, dimensional reduction
typically gives
\begin{equation}
\frac{1}{e_4^2}=\frac{2\pi R_\psi}{e_5^2}
\qquad\Longleftrightarrow\qquad
e_4^2=\frac{e_5^2}{2\pi R_\psi}.
\end{equation}
So compactification fixes the 4D coupling only if the parent coupling and radius are
already fixed by an independent principle in $S[\Theta]$.

\section{Decision on \texorpdfstring{$e^2$}{e2} status}
\textbf{Result:} At current proof state, UBT fixes the form of the U(1) kinetic operator,
but the absolute value of $e^2$ remains a normalization parameter.

\textbf{Therefore:} QED sector emergence is established, while full no-fit derivation of
$e^2$ from normalization alone is still open.

\end{document}
